\section{Locating the ceiling beams}\label{sec:beams}
Two complementary diagnostics localized the overhead structure. Autofocus fitted a thin layer to one view and scored its prediction of the other; the minimum gave $\AutofocusZ\pm\UncertaintyAutofocusZTotal\,\mathrm{m}$. Averaging shifted lattices reduced beam-pitch aliasing. Synthetic layer-centre recovery had signed bias $\ValidationFocusBias\,\mathrm{m}$ and maximum error $\ValidationFocusMaxError\,\mathrm{m}$.

Beam-shadow peaks supplied a geometric check. A feature at $(X_k,z)$ obeys
\begin{equation}
s_{x,p,k}=\frac{X_k-x_p}{z-z_p}.
\label{eq:triangulation}
\end{equation}
Matching transverse peaks between positions yielded $z_x=\BeamsZX\pm\UncertaintyBeamsZxTotal\,\mathrm{m}$. The joint-height solution gave beam pitch $\BeamsPitch\,\mathrm{m}$. This transverse estimate is emphasized because displacement along the beams is less identifiable. Autofocus measures an effective layer centre; triangulation localizes shadow features, so neither is automatically a bottom-face measurement.

\begin{figure}[htbp]
\centering
\includegraphics[width=\linewidth]{height.pdf}
\caption{One height diagnostic figure: cross-position autofocus (a) and transverse parallax (b). Dashed lines mark each estimator and shaded bands show its count-bootstrap spread.}\label{fig:height}
\end{figure}

Within the overlap of both views, the reconstructed layer at the triangulated height reproduced the raw-data beam positions with mean offset $\BeamsGateMeanAbsOffset\,\mathrm{m}$ and passed the verification gate. Single-view edge features were retained as separate diagnostics. This checks lateral placement, not beam depth or absolute scale.
